\section{Introduction}\label{sec:introduction}

Let $F$ be a cellular automaton on the full shift $A^{\Z^d}$ and let $\mu$ be
an input law.  Even when every finite input cylinder has positive mass,
iteration can redistribute mass in ways that make uniform output-cylinder
estimates difficult.  A concrete version of this tension appears in the
question of whether a surjective one-dimensional cellular automaton, started
from a full-support product law, can make the probability of a fixed cylinder
tend to zero~\cite[Question~1]{hellouinSaloTorma2026frequencies}.  The
arbitrary-surjective product case is not addressed here.  We instead isolate
two mechanisms that give a quantitative estimate at every time and on every
finite output shape.

Write $X=(X_v)_{v\in\Z^d}$ for the coordinate process with law $\mu$.  Our
probabilistic hypothesis is the following explicit all-other-sites
inequality: there is $\eta>0$ such that, for every $v\in\Z^d$ and $c\in A$,
\begin{equation}\label{eq:intro-sfe}
  \Prob\bigl(X_v=c\mid\cF_{\Z^d\setminus\{v\}}\bigr)\geq\eta
  \qquad\text{almost surely}.
\end{equation}
Here $\cF_C$ is generated by the coordinates in $C$.  No invariance,
independence, stationarity, mixing, Markov, or Gibbs assumption accompanies
\eqref{eq:intro-sfe}.  For a finite $S\subset\Z^d$ and $a\in A^S$, write
$[a]_S=\{x:x|_S=a\}$.

\Needspace{29\baselineskip}
\begin{theorem}\label{thm:main}
Let $d\geq1$, let $A$ be finite, and suppose that $\mu$ on $A^{\Z^d}$
satisfies \eqref{eq:intro-sfe}.  Let $F$ belong to one of the following two
classes.
\begin{enumerate}[label=\textup{(\roman*)}]
  \item $A=\F_q$ for a prime power $q$, and for some finite nonempty
  $N\subset\Z^d$,
  \[
      F(x)_v=b+\sum_{u\in N}c_u x_{v+u},
      \qquad b,c_u\in\F_q,
  \]
  where the coefficients $c_u$ are not all zero.
  \item For a finite nonempty $N\subset\Z^d$, the local rule
  $f\colon A^N\to A$ is permutive in a site $u_*\in N$, and a real linear
  functional $\ell\colon\mathbb R^d\to\mathbb R$ satisfies
  \[
      \ell(u_*)>\ell(u)\qquad(u\in N\setminus\{u_*\}).
  \]
\end{enumerate}
Then, for every integer $t\geq0$, every finite $S\subset\Z^d$, and every
$a\in A^S$,
\begin{equation}\label{eq:main-bound}
    (F^t_*\mu)([a]_S)\geq\eta^{|S|}.
\end{equation}
For $S=\varnothing$, \eqref{eq:main-bound} is the equality $1=\eta^0$.
\end{theorem}

The common probabilistic engine is a finite conditional block estimate.
Condition \eqref{eq:intro-sfe} pays one factor of $\eta$ for each site in a
specified block, even when the desired block pattern is a measurable
function of all sites outside the block.  Thus the dynamical task is to find
exactly $|S|$ whole-site pivots whose values are uniquely determined by the
target output pattern after every other input site is fixed.

For the scalar affine class, let
\[
  P(z)=\sum_{u\in N}c_u z^u
  \in\F_q[z_1^{\pm1},\ldots,z_d^{\pm1}].
\]
The time-$t$ linear kernel is $P^t$.  A dependence among the output rows
indexed by $S$ would give $A(z)P(z)^t=0$ for a nonzero finite Laurent
polynomial $A$.  The Laurent ring is an integral domain, so all rows are
independent and an invertible set of $|S|$ site columns can be selected.

For the exposed-corner class, $tu_*$ is the unique maximizer of $\ell$ on the
$t$-fold Minkowski sum $tN$, and the time-$t$ rule remains permutive in the
raw coordinate $tu_*$.  The candidate pivots are $s+tu_*$ for $s\in S$.
If the pivot belonging to $s$ influences the output at a distinct $r$, then
$\ell(s)<\ell(r)$.  Increasing $\ell$-order therefore turns the finite
pivot map into a triangular bijection, including when $S$ is disconnected or
nonconvex.

\FloatBarrier
\begin{table}[t]
  \centering
  \small
  \setlength{\tabcolsep}{3.4pt}
  \renewcommand{\arraystretch}{1.13}
  \caption{Mechanisms and scope.  Only the two upper rows are theorem
  regimes.}
  \label{tab:mechanisms}
  \begin{tabularx}{\textwidth}{@{}
    >{\RaggedRight\arraybackslash}p{25mm}
    >{\RaggedRight\arraybackslash}p{38mm}
    >{\RaggedRight\arraybackslash}X
    >{\RaggedRight\arraybackslash}p{31mm}@{}}
    \toprule
    Regime & Load-bearing hypothesis & Whole-site pivot mechanism
      & Status \\
    \midrule
    Scalar affine
      & Nonzero Laurent polynomial over the field $\F_q$
      & Full row rank gives an invertible $|S|$-site column minor
      & Bound $\eta^{|S|}$ for all $t,S$ \\
    Exposed corner
      & One permutive site uniquely exposed by a real functional
      & Translated corners form a triangular bijection in functional order
      & Bound $\eta^{|S|}$ for all $t,S$ \\
    Vector affine
      & Whole-site energy plus scalar-coordinate rank only
      & Scalar pivots may occupy more than $|S|$ input sites
      & Blanket bound false \\
    General surjective
      & Surjectivity without an exposed permutive coordinate
      & No triangular pivot follows from the hypothesis
      & Not claimed \\
    \bottomrule
  \end{tabularx}
\end{table}


\FloatBarrier

Affine and Abelian cellular automata have an extensive asymptotic
randomization theory.  Pivato and Yassawi obtain Haar convergence in density
under diffusivity and harmonic-mixing hypotheses
\cite[Theorem~12 and Proposition~14]{pivatoYassawi2002limit}
\cite[Theorem~3, Proposition~12, and Theorem~15]{pivatoYassawi2004limitII},
while later work
characterizes Ces\`aro randomization for Abelian rules through the absence of
solitons~\cite[Theorem~2]{hellouinSaloTheyssier2020randomization}.  The estimate
\eqref{eq:main-bound} has a different quantifier pattern: it is a lower bound
at each time and permits dependent laws under \eqref{eq:intro-sfe}; it is not
a convergence assertion.

Permutivity, corner-filling, extremal geometry, and order-based extension
arguments occur in the theory of TEP subshifts
\cite[Theorems~1, 19, and 20]{salo2022cutting}.  Our exposed-corner proof uses that
geometry only to build a finite-time pivot bijection.  Work on statistical
mechanics of surjective cellular automata studies invariant and equilibrium
measures under different assumptions
\cite[Corollary~21 and Theorem~41]{kariTaati2015statistical}; the
term \emph{Gibbs} by itself is not a substitute for
\eqref{eq:intro-sfe}.

Two boundaries are part of the statement rather than optional
generalizations.  Scalar-coordinate rank does not imply the same estimate on
a vector alphabet when $\eta$ is a whole-site conditional constant; a
reversible counterexample is given in \cref{prop:vector-firewall}.
Surjectivity alone also does not imply the exposed-permutive hypothesis;
\cref{prop:surjective-firewall} gives a finite exact certificate for an onto
binary rule with no permutive coordinate.  That rule is used only to separate
hypotheses, not as a counterexample to \eqref{eq:main-bound}.

The rest of the note is organized as follows.
\Cref{sec:finite-energy} proves the conditional block lemma and the measurable
pivot principle.  \Cref{sec:affine} proves the scalar affine case.
\Cref{sec:corner} proves persistence of an exposed corner and the
arbitrary-shape triangular solve.  \Cref{sec:scope} proves sharpness, records
the two firewalls, and describes exact finite controls.
